\section{Introduction}
\label{sec:introduction}

Modern retrieval systems increasingly rely on both dense and sparse representations. Dense embeddings encode semantic similarity and are effective for paraphrases, concept-level matching, and natural-language queries. Sparse vectors preserve lexical evidence and are important for exact terms, entity names, rare words, and identifier-like queries. These two signals are complementary: dense retrieval improves semantic recall, while sparse retrieval protects lexical precision. Dense--sparse hybrid retrieval has therefore become a common design in search engines, question answering systems, recommendation, and retrieval-augmented generation.

The dominant architecture for hybrid retrieval is a two-route pipeline. A dense approximate nearest neighbor (ANN) index retrieves candidates by dense distance, and a sparse index retrieves candidates by sparse distance. The two candidate sets are then merged and reranked with a hybrid distance, usually a weighted combination of dense and sparse distances. This design is attractive because it reuses mature dense and sparse indexes. However, it treats hybrid retrieval as a late-fusion problem: the indexes decide which objects are even considered, and the hybrid distance is applied only after each single-modality route has already discarded most objects.

This late-fusion design can miss important hybrid neighbors. A relevant object may have a moderately small dense distance and a moderately small sparse distance, but become highly ranked after the two distances are combined. Such an object may not appear in either the dense top candidates or the sparse top candidates, so reranking cannot recover it. Increasing both candidate budgets can reduce this failure case, but it also increases latency, duplicates search effort, and introduces additional tuning knobs. More importantly, larger candidate sets still do not create an index topology aligned with the hybrid distance.

Graph-based ANN search makes this topology issue central. A graph index is effective when its edges connect objects that are close under the same metric used at query time. A dense graph is navigable under dense distance, and a sparse graph is navigable under sparse distance. Hybrid retrieval, however, ranks objects by a combined dense--sparse distance. Searching a single-modality graph and reranking visited nodes therefore still uses a graph built for the wrong objective. A natural alternative is to build a unified graph over hybrid objects and search it directly with the hybrid distance.

Building such a graph is non-trivial because hybrid retrieval is controlled by a weight $\alpha$ that changes the relative importance of dense and sparse distances. A graph constructed for one fixed weight may preserve the right neighbors for that weight but miss neighbors needed by sparse-dominant or dense-dominant queries. Rebuilding a separate graph for every possible $\alpha$ is impractical because graph construction is expensive and the desired weight may vary across datasets, applications, or query types. A practical hybrid graph should therefore be hybrid-aware, robust to different $\alpha$ values, and still sparse enough for efficient graph traversal.

To address these limitations, we move hybrid awareness from the reranking stage into the graph construction stage. We propose \textbf{Unified Hybrid Graph (UHG)}, a graph index for dense--sparse hybrid retrieval under varying hybrid weights. UHG first obtains candidate neighbors from both dense and sparse evidence. It then evaluates these candidates under a grid of hybrid weights and retains candidates that become nearest under any sampled weight. We call this procedure alpha-cover neighbor selection. The resulting directed graph is further refined with reverse-edge augmentation, relative-neighborhood-inspired pruning, and connectivity repair, producing a bounded-degree graph for hybrid ANN search. At query time, UHG traverses this graph using the query-time hybrid distance.

We also study \textbf{Unified Hybrid Graph with SINDI (UHGS)}, a sparse-guided search variant of UHG. Instead of using the sparse index as an independent retrieval route whose results are merely merged with dense results, UHGS uses sparse retrieval results as query-dependent entry points for UHG traversal. It further reuses the sparse distances computed by SINDI during its search, attaching them to the hybrid query so that UHG can avoid redundant sparse scoring and prune unnecessary dense distance computations. This design turns sparse retrieval into both a routing and pruning signal for hybrid graph search, allowing dense and sparse evidence to interact during traversal rather than only during final reranking.

This paper makes the following contributions:
\begin{itemize}[leftmargin=*]
    \item We formulate the topology mismatch in dense--sparse hybrid retrieval: existing two-route pipelines use a hybrid distance for final ranking but rely on separate single-modality structures for candidate generation.
    \item We propose UHG, a unified hybrid graph index that selects alpha-cover neighbors across a family of dense--sparse hybrid weights, reducing the need to rebuild the graph for different $\alpha$ values.
    % \item We refine the alpha-cover graph with reverse-edge augmentation, relative-neighborhood-inspired pruning, and connectivity repair so that the graph remains bounded-degree and navigable.
    \item We propose UHGS, which uses sparse retrieval results as entry points for hybrid graph traversal and studies sparse-guided routing as an alternative to two-route candidate fusion.
    \item We reuse the sparse distances computed by SINDI during hybrid graph search and introduce a sparse-distance-based pruning strategy that reduces dense distance computations at query time.
    \item We design an evaluation framework comparing UHG and UHGS with two-route hybrid baselines, including HGraph+HGraph and HGraph+SINDI, using recall and QPS.
\end{itemize}

The remainder of this paper is organized as follows. Section~\ref{sec:problemDefinition} reviews ANNS, graph-based ANNS, and dense--sparse hybrid search. Section~\ref{sec:solution1} motivates unified hybrid graph construction. Section~\ref{sec:solution2} presents UHG and UHGS. Section~\ref{sec:experimental} describes the experimental design, and Section~\ref{sec:conclusion} concludes the paper.
